\section{Discussion}
\label{sec:discussion}

The comparison among research interfaces suggests that mathematical capability becomes useful for autonomous co-design when it is organized around a task-specific physical question. Direct Mathematical Tools provides numerical analysis and optimization, while \softrobotagent{} couples mathematical execution to explicit physical hypotheses, observables, model bindings, and control treatments. This contrast highlights research-question construction as a central organizing capability: the agent identifies the physical relationship underlying a task bottleneck and selects an analysis whose response matches that relationship. The resulting evidence then informs the next structural or control intervention.

The representative cases reveal complementary relationships among morphology, deformation, and actuation. In Passive Tip Suppression, tendon routing shapes the admissible deformation directions, while material and cross-sectional properties govern their load-dependent response. In Shared Morphology, the effect of a structural change depends on the control treatment: a revised morphology can reduce tendon-speed demand under fixed input while increasing position-holding error, whereas control adaptation can recover task feasibility on the same structure. Figure-eight Tracking illustrates how equilibrium analysis connects configuration-dependent tendon forces to trajectory execution through a continuous feedforward schedule. Together, these findings show that structural quality is task-conditioned and depends on both the physical response of the morphology and the control strategy used to exploit it.

Research continuity preserves two complementary forms of accumulated experience. Knowledge retains physical relationships, design rationales, and their applicability conditions, whereas Skill retains reusable methods for investigating and validating physical hypotheses. Their complementary contributions indicate that continued research benefits from preserving both what has been established about the physical system and how that knowledge was obtained. Reporter interpretation and hierarchical retrieval further show that evidence summaries and access to original records serve distinct roles: summaries support efficient interpretation, while source-level retrieval preserves access to decision-critical evidence. Research memory therefore benefits from combining concise synthesis with traceable access to the evidence underlying each finding.

Task-conditioned mathematical research provides an organizing principle for mechanism-guided structure--control co-design. Extending this process to richer material dynamics, broader tendon-routing families, and more complex loading and contact conditions would expand the physical relationships available for investigation. Fabricated-robot experiments offer the next step toward connecting model-based sensitivities and predicted responses with measured behavior, allowing accumulated research evidence to incorporate physical measurements and model applicability conditions.
